% Generated draft: requires mathematical and source review.
\documentclass[11pt]{article}
\usepackage[margin=25mm]{geometry}
\usepackage{fontspec}
\setmainfont{Latin Modern Roman}
\usepackage{amsmath,amssymb,mathtools,unicode-math}
\setmathfont{Latin Modern Math}
\usepackage{enumitem,needspace,xurl,hyperref}
\setlength{\parindent}{0pt}
\setlength{\parskip}{5pt}
\setlength{\emergencystretch}{3em}
\newcommand{\N}{\mathbb{N}}
\newcommand{\Z}{\mathbb{Z}}
\newcommand{\Q}{\mathbb{Q}}
\newcommand{\R}{\mathbb{R}}
\newcommand{\C}{\mathbb{C}}
\newcommand{\F}{\mathbb{F}}
\newcommand{\A}{\mathbb{A}}
\newcommand{\PP}{\mathbb{P}}
\newcommand{\E}{\mathbb{E}}
\newcommand{\eps}{\varepsilon}
\newcommand{\dd}{\,\mathrm{d}}
\begin{document}
\section*{UnsolvedMath ID: 700241}
\textbf{Source problem code:} catalogue entry
\subsection{Definitions and mathematical statement}
For a group \(G\), define \(G^{(0)}=G\) and \(G^{(i+1)}=[G^{(i)},G^{(i)}]\). Under the standard interpretation of the variety in the source, for a positive integer \(n\) let
\[
\mathcal S_n=\{G:G^{(n)}=1\}.
\]
Thus the derived length is at most \(n\); exact derived length \(n\) would require the additional condition \(G^{(n-1)}\ne1\), which the source does not explicitly impose.

For a set \(X\), let \(F(X)\) be the free group on \(X\), and put \(S_n(X)=F(X)/F(X)^{(n)}\). Interpret the source's class \(\mathfrak K_n\) as the groups admitting a presentation with one defining relation relative to \(\mathcal S_n\):
\[
G\cong S_n(X)/\langle\!\langle r\rangle\!\rangle_{S_n(X)},\qquad r\in S_n(X),
\]
where \(\langle\!\langle r\rangle\!\rangle_{S_n(X)}\) is the smallest normal subgroup of \(S_n(X)\) containing \(r\). The source does not specify whether \(X\) must be finite or whether a trivial or redundant defining relation is permitted; these conventions require review.

\begin{enumerate}
\item[(a)] Determine necessary and sufficient conditions for a group \(G\in\mathfrak K_n\) to have nontrivial centre, where
\[
Z(G)=\{z\in G:zg=gz\text{ for every }g\in G\}.
\]
In particular, for \(n\ge2\), determine whether there exists \(G\in\mathfrak K_n\) such that
\[
Z(G)\ne\{1\}
\quad\text{and}\quad
\text{there are no }x,y\in G\text{ with }G=\langle x,y\rangle.
\]
The question is posed with \(n\) as a parameter. The supplied editors' comment, dated 1998, reports that the questions in part (a) were answered for \(n=2\); it does not supply those answers.
\item[(b)] Describe the abelian subgroups \(A\le G\) of groups \(G\in\mathfrak K_n\), where abelian means \([a,b]=1\) for every \(a,b\in A\).
\item[(c)] Investigate the periodic subgroups \(P\le G\) of groups \(G\in\mathfrak K_n\), where periodic means
\[
\forall p\in P\ \exists m\in\Z_{\ge1}\quad p^m=1.
\]
No finiteness assumption on these subgroups is stated.
\end{enumerate}
Parts (b) and (c) are structural research requests: the source specifies neither a proposed classification nor a particular assertion to prove.
\subsection{Short English statement}
M. I. Kargapolov asks when a group with one defining relation in a soluble variety has nontrivial centre, especially when it cannot be generated by two elements, and asks for descriptions of its abelian and periodic subgroups. The supplied 1998 editors' comment reports answers to the centre questions for \(n=2\). Presentation conventions and the intended scope of the subgroup descriptions require human review.
\subsection{Sources}
\begin{itemize}
\item Local UnsolvedMath registry, supplied complete record with permanent ID \(700241\), exact source code \texttt{catalogue entry}, title \emph{Kourovka Notebook Problem 4.33}. The record labels the problem open; this is supplied registry metadata, not an independently verified assessment.
\item M. I. Kargapolov, problem 4.33, \emph{The Kourovka Notebook (Unsolved Problems in Group Theory)}, 4th issue (1973). The supplied record identifies the later source as the 21st issue edition (2026), \texttt{arXiv:1401.0300v48}. The collection description names E. I. Khukhro and V. D. Mazurov as editors.
\item Editors' comment (1998), supplied within part (a): the centre questions were answered for \(n=2\), citing E. I. Timoshenko, \emph{Siberian Math. J.}, \textbf{14}, no. 6 (1973), 954--957; and \emph{Math. Notes}, \textbf{64}, no. 6 (1998), 798--803. Article titles and links are not supplied. This dated comment is reproduced as source attribution and has not been independently verified.
\end{itemize}
\end{document}
